\chapter{Reasoning Violation Taxonomy in Full Detail}
\label{app:taxonomy}

This appendix reproduces the V1--V5 taxonomy in full, with the rule-based
detection logic, the lexicon term lists, and the safe-negation patterns that
the production checker uses. The canonical source of truth is
\texttt{src/reason\_guard\_checker.py}.

\section{V1: Fabricated Reasoning}
\label{sec:appendix-v1}

\paragraph{Definition.}
The AI asserts facts, causes, or entities not present in or derivable from
the formal bound. Under the proxy bound, every \texttt{*\_claim\_allowed}
flag in \texttt{machine\_constraints} defaults to \texttt{false}; therefore
any unsupported attack, mitigation, causal, or numeric assertion fires V1.

\paragraph{Detection.}
For each candidate term, the detector applies a two-layer check:

\begin{enumerate}
  \item A spaCy dependency-parse filter
        (\texttt{claim\_extraction/negation\_scope.py}) walks from the term
        token up the dependency tree and treats the term as safe if any
        ancestor has a \texttt{neg} child or a child whose lowered text is
        in the negation set \{\texttt{not, no, never, neither, nor,
        without}\}. The term is also treated as safe when its sentence
        contains \texttt{proxy} or \texttt{schema}.
  \item A substring-based safe-negation list
        (\texttt{SAFE\_NEGATION\_PATTERNS}) provides the same protection
        when spaCy is unavailable. Both layers are intentionally
        conservative: a term is treated as unsafe whenever neither layer
        can match it to a negation.
\end{enumerate}

\paragraph{Lexicon (excerpts).}
The full lists are defined as Python constants in
\texttt{reason\_guard\_checker.py}.

\begin{itemize}
  \item Attack terms: attack, malicious, replay, flooding, mitm,
        man in the middle, scan, scanning, dos, denial of service,
        command injection, intrusion, compromise, attacker, exploit.
  \item Mitigation terms: isolate, isolated, shutdown, block, blocked,
        quarantine, disconnect, emergency response, immediate response,
        prevent further.
  \item Causal terms: because, caused by, due to, therefore, in order to,
        as a result, consequently, thus the cause, the cause of.
  \item Numeric check: each integer in the response is compared against the
        bound's feature values plus the constants $\{104, 2404\}$ plus the
        observed timestamp; values not in this set and greater than 1\,000
        are flagged.
\end{itemize}

\section{V2: Contradicted Reasoning}
\label{sec:appendix-v2}

\paragraph{Definition.}
The AI explicitly contradicts a fact established by the formal bound.

\paragraph{Detection.}
The detector currently checks two contradiction cases:

\begin{enumerate}
  \item Direction inversion: the response uses the opposite of the bound's
        \texttt{communication\_direction} and does not also mention the
        bound's direction.
  \item Identifier substitution: the response contains the literal
        \texttt{destination\_id} value alongside the keyword \emph{source}
        without containing the literal \texttt{source\_id} value.
\end{enumerate}

Additional contradiction cases will be added as they are observed in the
M4 annotation pass.

\section{V3: Over-Generalised Reasoning}
\label{sec:appendix-v3}

\paragraph{Definition.}
The AI produces a correct high-level category but loses formally
established specifics that the operator needs in order to act.

\paragraph{Detection.}
The detector fires V3 when any normality term appears in the response
together with fewer than three of the formally observed specific values
(source\_id, destination\_id, srcport, dstport, asdu\_type, asdu\_cot).
Normality terms: normal traffic, benign, routine, typical, standard
communication, no specific concern, normality.

\section{V4: Under-Specified Reasoning}
\label{sec:appendix-v4}

\paragraph{Definition.}
The AI omits formally established information that is operationally
required for correct response.

\paragraph{Detection.}
The detector fires V4 when any of the following required references is
missing from the response: source, destination, port, ASDU, and the
combined uncertainty / proxy clarification. Missing references are
reported as individual reason strings so the operator can correct the
explanation rather than only learn that V4 has fired.

\section{V5: Incoherent Reasoning}
\label{sec:appendix-v5}

\paragraph{Definition.}
The AI explanation is internally self-contradictory, independent of the
formal output.

\paragraph{Detection.}
The detector uses spaCy sentence segmentation
(\texttt{claim\_extraction/coherence.py}) to find sentences that contain
both members of a canonical conflicting pair. The detector explicitly
skips sentences that contain a negation marker so that statements such as
\emph{the bound does not confirm maliciousness or normality} are not
treated as incoherent. The substring-based fallback applies when spaCy is
unavailable.

\paragraph{Conflicting pairs.}
\begin{itemize}
  \item \emph{normal} and \emph{malicious}
  \item \emph{normal} and \emph{attack}
  \item \emph{benign} and \emph{malicious}
  \item \emph{routine} and \emph{attack}
  \item \emph{isolated} and \emph{entire network}
  \item \emph{single device} and \emph{entire network}
  \item \emph{client-to-server} and \emph{server-to-client}
\end{itemize}

\section{Severity Mapping}
\label{sec:appendix-severity}

The mapping from the violation set to the severity level is:

\begin{center}
\begin{tabular}{ll}
  \toprule
  Violation set & Severity \\
  \midrule
  $\{$ V1 $\}$, $\{$ V2 $\}$, $\{$ V5 $\}$ (or any superset) & high \\
  $\{$ V3, V4 $\}$ (without any high-severity class) & medium \\
  $\{$ V3 $\}$ or $\{$ V4 $\}$ alone & low \\
  $\emptyset$ & none \\
  \bottomrule
\end{tabular}
\end{center}

The mapping intentionally favours operator caution: any indication of
fabricated, contradicted, or incoherent reasoning escalates the
explanation to the high-severity review queue described in
Section~\ref{sec:discussion-deployment}.
